\subsection{Особенности блокчейна Bitcoin как источника данных}

Блокчейн Bitcoin является распределенной системой, в которой данные о
транзакциях публикуются и подтверждаются в рамках механизма консенсуса.
Практический доступ к этим данным в проектируемой системе осуществляется через
узел Bitcoin Core по \gls{jsonrpc}. Такой способ взаимодействия обеспечивает
достоверность источника, но при этом ориентирован прежде всего на обслуживание
сетевого узла и администрирование, а не на выполнение аналитических или
массовых прикладных выборок \cite{btc_core_rpc}.

В качестве источника данных блокчейн Bitcoin характеризуется следующими
особенностями:
\begin{enumerate}
    \item блочная организация данных, при которой транзакции доступны в составе
    последовательности блоков, связанных отношением предшествования;
    \item модель \gls{utxo}, требующая восстановления состояния адресов на
    основе истории создания и расходования выходов;
    \item наличие mempool как отдельного множества неподтвержденных транзакций,
    состояние которого изменяется независимо от подтвержденной цепи;
    \item возможность reorg, то есть пересмотра канонической ветки на
    ограниченной глубине;
    \item значительный объем исторических данных, требующий пакетной и
    устойчивой к повторной обработке индексации.
\end{enumerate}

С инженерной точки зрения это означает, что прикладная система не может
ограничиться единичными RPC-вызовами к узлу. Для получения данных,
пригодных для дальнейшего использования, необходимо организовать полный цикл:
получение блока, декомпозицию его содержимого, запись транзакций, входов и
выходов в реляционную модель, обновление адресных индексов, формирование
исторических агрегатов и синхронизацию с текущим состоянием mempool.

В проекте \texttt{Bitcoin Blockchain Indexer} эти особенности учтены через
разделение контура доступа к узлу и контура прикладного хранения данных.
Bitcoin Core используется как доверенный источник первичной информации, тогда
как проектируемая система обеспечивает нормализацию данных и подготовку
устойчивых прикладных представлений для последующей выдачи пользователям и
внешним сервисам.
